\documentclass[10pt,aspectratio=169,compress]{beamer}
\usepackage[english]{babel}
\IfFileExists{kotex.sty}{\usepackage{kotex}\newif\ifABKorean\ABKoreantrue}{\newif\ifABKorean\ABKoreanfalse}
\usepackage{amsmath,amssymb}
\IfFileExists{newtxtext.sty}{\usepackage{newtxtext,newtxmath}}{}
\usepackage{theme/academic-beamer}
\usepackage{pgfplots}
\pgfplotsset{compat=1.18}
% Deck-specific settings. No school logo or research identity is bundled.
\AcademicPreset{classic}
% Optional deck-local overrides go AFTER the preset:
% \LayoutSetup{columns=8,gutter=4mm}
% \AcademicPalette{neutral}
% Choose none, start, sections, or both. The demo shows a single opening outline.
\AcademicOutlineSetup{mode=start,title=Outline,toc-options=hideallsubsections}
\AcademicIdentity{Sample University}
\title[Evidence and local outcomes]{Evidence and Local Outcomes}
\subtitle{An illustrative academic presentation}
\author[Author A \& B]{\texorpdfstring{Author A \qquad Author B}{Author A and Author B}}
\date[Seminar]{Research Seminar}
% \AcademicLogo{assets/user-provided-logo.pdf}

\hypersetup{pdftitle={Academic Beamer layout examples},pdfauthor={Illustrative authors}}
\begin{document}
\AcademicTitleFrame
\AcademicOutline
\section[Question]{Research Question}
\begin{frame}{Research Question}
  \begin{GridRow}
    \GridBlank{1}
    \begin{GridCell}{4}
      \vspace{9mm}\centering
      {\fontsize{16}{20}\selectfont How does a large income payment reach local markets?\par}
      \vspace{5mm}
      \ifABKorean 연구질문을 한 문장으로 전달하고, 여백을 남깁니다.\else
      A focused question can stand on its own.\fi
    \end{GridCell}
    \GridBlank{1}
  \end{GridRow}
\end{frame}
\section[Data]{Data}
\begin{frame}{Data Sources}
  \begin{GridRow}
    \begin{GridCell}{2}
      \begin{block}{Movement}
        \begin{itemize}\setlength\itemsep{1mm}
          \item \ifABKorean 근로자의 거주지역\else Residential origins\fi
          \item \ifABKorean 직장과 통근 연결\else Workplace links\fi
        \end{itemize}
      \end{block}
    \end{GridCell}
    \begin{GridCell}{2}
      \begin{block}{Local sales}
        \begin{itemize}\setlength\itemsep{1mm}
          \item \ifABKorean 가맹점 소재지 매출\else Merchant location\fi
          \item \ifABKorean 월별 소비 변화\else Monthly outcomes\fi
        \end{itemize}
      \end{block}
    \end{GridCell}
    \begin{GridCell}{2}
      \begin{block}{Background}
        \begin{itemize}\setlength\itemsep{1mm}
          \item \ifABKorean 지역 인구 구성\else Population\fi
          \item \ifABKorean 주거환경과 특성\else Housing conditions\fi
        \end{itemize}
      \end{block}
    \end{GridCell}
  \end{GridRow}
  \ifABKorean\small 세 자료의 관측 단위와 역할을 구분합니다.\else
  \small Distinguish what each source actually measures.\fi
  \AcademicNote{These sources and all values in this deck are illustrative.}
\end{frame}
\begin{frame}{Summary Statistics}
  \begin{GridRow}
    \begin{GridCell}{6}
      \begin{AcademicStatTable}{Xrrr}
        \textbf{Variable} & \textbf{Mean} & \textbf{SD} & \textbf{Units}\\\midrule
        \multicolumn{4}{l}{\textbf{Panel A: Exposure and population}}\\\addlinespace[1mm]
        Exposure & 0.32 & 0.74 & Percent\\
        Population & 18.4 & 9.7 & Thousands\\
        \addlinespace[2mm]
        \multicolumn{4}{l}{\textbf{Panel B: Outcome}}\\\addlinespace[1mm]
        Monthly sales & 4.8 & 2.1 & Index units\\
      \end{AcademicStatTable}
    \end{GridCell}
  \end{GridRow}
  \begin{itemize}\setlength\itemsep{2mm}
    \item A common sample connects exposure and outcomes.
    \item Summary statistics follow the data definitions.
  \end{itemize}
  \AcademicNote{Synthetic example only. These numbers are not research findings.}
\end{frame}
\begin{frame}{Measurement Concepts}
  \begin{GridRow}
    \begin{GridCell}{6}
      \begin{AcademicRelationTable}{>{\raggedright\arraybackslash}X >{\raggedright\arraybackslash}X >{\raggedright\arraybackslash}X}
        \AcademicTableHeading{Concept} & \AcademicTableHeading{What is observed} & \AcademicTableHeading{Interpretation}\\[1mm]
        \AcademicTableLabel{Household spending} &
          \ifABKorean 가구가 지출한 소비\else Purchases by a household\fi &
          \ifABKorean 개인 소비 반응\else Household consumption response\fi\\[2mm]
        \AcademicTableLabel{Local merchant sales} &
          \ifABKorean 지역 가맹점에서 발생한 매출\else Sales at local merchants\fi &
          \ifABKorean 지역 상권에 남은 소비\else Local-market footprint\fi\\[2mm]
        \AcademicTableLabel{Vehicle registration} &
          \ifABKorean 거주지 기준 신규등록\else Registrations by residence\fi &
          \ifABKorean 내구재 구매의 보완 지표\else Complementary durable-goods measure\fi\\
      \end{AcademicRelationTable}
    \end{GridCell}
  \end{GridRow}
  \AcademicNote{Illustrative comparison. Accent labels and whitespace organize concepts; no rules or background fills.}
\end{frame}
\section[Measurement]{Measurement}
\begin{frame}{Income and the Local-Sales Footprint}
  \begin{GridRow}
    \begin{GridCell}{6}
      \begin{block}{An accounting framework}
        \AcademicEquation{\Delta Y_{rh}=q_{rh}\lambda_{rh}m_{rh}F_r+G_{rh}}
        \ifABKorean
          $F_r$: 지역에 배분된 소득, $G_{rh}$: 기타 지역 변화\\[2mm]
          $m$: 소비 전환, $\lambda$: 지역 내 지출, $q$: 자료 포착 비중
        \else
          $F_r$: allocated income; $G_{rh}$: other local changes.\\[2mm]
          $m$: consumption share; $\lambda$: local share; $q$: observed share.
        \fi
      \end{block}
    \end{GridCell}
  \end{GridRow}
  \small An accounting identity does not identify each behavioral parameter.
\end{frame}
\section[Evidence]{Evidence}
\begin{frame}{An Illustrative Outcome Path}
  \begin{GridRow}
    \begin{GridCell}{4}
      \begin{AcademicFigurePanel}{Observed and Comparison Paths}{Illustrative monthly index; both paths are code-generated.}
      \centering
      \begin{tikzpicture}
        \begin{axis}[width=\linewidth,height=40mm,scale only axis=false,
          xmin=1,xmax=8,ymin=90,ymax=112,
          xlabel={Month},ylabel={Index},
          tick label style={font=\footnotesize},label style={font=\small},
          legend style={font=\scriptsize,draw=none,at={(0.02,0.98)},anchor=north west},
          axis lines=left,grid=major,grid style={gray!15}]
          \addplot[ABInk,thick]coordinates{(1,96)(2,98)(3,100)(4,102)(5,108)(6,106)(7,102)(8,101)};
          \addlegendentry{Observed}
          \addplot[blue!65!black,thick,dashed]coordinates{(1,97)(2,98)(3,100)(4,101)(5,103)(6,102)(7,101)(8,100)};
          \addlegendentry{Comparison}
        \end{axis}
      \end{tikzpicture}
      \end{AcademicFigurePanel}
    \end{GridCell}
    \begin{GridCell}{2}
      \begin{alertblock}{Timing}
        \ifABKorean 지급 이전과 이후의 경로를 함께 봅니다.\else
          Compare the path before and after payment.\fi
      \end{alertblock}
      \vspace{2mm}\small
      \ifABKorean 관측과 가설을 구분합니다.\else
      Separate observations from explanations.\fi
    \end{GridCell}
  \end{GridRow}
  \AcademicNote{Illustrative data only; no empirical or causal claim.}
\end{frame}
\begin{frame}{Conclusion}
  \begin{GridRow}
    \begin{GridCell}{3}
      \begin{alertblock}{Answer}
        Return to the original question and summarize the supported finding.
      \end{alertblock}
    \end{GridCell}
    \begin{GridCell}{3}
      \begin{block}{Possible mechanism}
        State an explanation as a hypothesis unless the study identifies it.
      \end{block}
    \end{GridCell}
  \end{GridRow}
  \AcademicNote{Use only conclusions supported by the actual study.}
\end{frame}
\begin{frame}\centering\LARGE Thank you\\[5mm]\large Q\&A\end{frame}
\AcademicAppendix
\begin{frame}{Appendix Outline}
  \begin{itemize}\setlength\itemsep{3mm}
    \item Data and definitions
    \item Identification diagnostics
    \item Additional results and references
  \end{itemize}
\end{frame}
\begin{frame}{An Illustrative Estimation Table}
  \begin{GridRow}
    \begin{GridCell}{6}
      \begin{AcademicStatTable}{Xrr}
        & \textbf{(1)} & \textbf{(2)}\\
        & Baseline & Alternative\\\midrule
        Exposure coefficient & 0.12 & 0.08\\
        & (0.09) & (0.11)\\\addlinespace[2mm]
        Observations & 1,200 & 1,180\\
        Fixed effects & Yes & Yes\\
      \end{AcademicStatTable}
    \end{GridCell}
  \end{GridRow}
  \AcademicNote{Synthetic example only. Parentheses illustrate standard errors; no empirical inference or significance is implied.}
\end{frame}
\end{document}
